% Two-page committee summary, 23 September 2026.
% Written to be read cold, without the full progress report and without
% familiarity with the project. Self-contained: compiles on its own.
\documentclass[11pt,a4paper]{article}

\usepackage[T1]{fontenc}
\usepackage[utf8]{inputenc}
\usepackage{lmodern}
\usepackage[margin=22mm]{geometry}
\usepackage{microtype}
\usepackage{booktabs}
\usepackage{tabularx}
\usepackage{siunitx}
\usepackage{enumitem}
\usepackage{xcolor}
\usepackage{titlesec}

\setlist{nosep}
\titlespacing*{\section}{0pt}{1.4ex plus .2ex}{0.6ex}
\titleformat{\section}{\normalfont\large\bfseries}{\thesection.}{0.6em}{}
\pagestyle{empty}

\begin{document}

\begin{center}
  {\Large\bfseries Vision-Based Navigation for AMRs in a Dynamic Warehouse}\\[0.4em]
  {\large Progress summary for committee review}\\[0.3em]
  Bolton Athit Davies \quad$\cdot$\quad 23 September 2026 \quad$\cdot$\quad
  evidence cutoff 16 September 2026
\end{center}

\vspace{0.6em}
\noindent\rule{\textwidth}{0.4pt}
\vspace{0.4em}

\section{The question}

Can a vision-based pipeline provide reliable localization for an autonomous
mobile robot in a warehouse that contains moving objects? The work compares two
stereo-inertial front ends, ORB-SLAM3 and VINS-Fusion, under controlled changes
in scene dynamics, floor texture and compute load, and asks whether masking
dynamic objects with a learned detector improves localization. Path planning is
outside the scope: the deliverable is a reliable navigation state and map, not a
planner.

\section{What exists}

A warehouse simulator with static, dynamic, textured-floor and featureless-floor
variants; both estimators instrumented to record per-frame timing, delivery,
state, CPU and memory; a custom YOLO detector feeding dynamic-object masks into
both; RTAB-Map as a separate mapping and loop-closure back end; and a
run-logging and validity-audit harness. The evidence base is 66 logged estimator
runs across seven simulation datasets at three repeats per condition, five
multi-sensor recordings spanning four worlds, \num{9860} camera frames with
projected object ground truth, and one reconstructed map with geometry and
drivability evaluation.

\section{What is established}

\textbf{A wheel+IMU EKF currently outperforms both visual pipelines in static
scenes.} On five recordings spanning four worlds, the EKF achieves the lowest
aligned ATE on four of five, at \SIrange{0.10}{0.88}{\percent} drift.

\begin{center}\small
\begin{tabular}{lrrrrr}
\toprule
Dataset & IMU only & Wheel odom. & VINS-Fusion & ORB-SLAM3 & EKF \\
\midrule
\texttt{allsensor\_000} & 21.70 & 2.80 & 0.62 & 0.58 & \textbf{0.19} \\
\texttt{allsensor\_001} & 12.62 & 0.98 & 0.38 & 0.37 & \textbf{0.31} \\
\texttt{allsensor\_002} & 13.24 & 1.49 & 0.56 & 0.60 & \textbf{0.26} \\
\texttt{allsensor\_003} & 45.56 & 3.19 & 0.78 & 6.48 & \textbf{0.42} \\
\texttt{allsensor\_004} & 26.71 & 3.24 & \textbf{0.43} & 1.24 & 0.49 \\
\bottomrule
\end{tabular}\\[0.3em]
\footnotesize Aligned ATE RMSE in metres, rigid $SE(3)$ fit with scale fixed.
\end{center}

This is the baseline the project previously lacked: until now every trajectory
figure reported visual error with nothing to compare it against. The EKF's
advantage is specific to these conditions --- static scene, simulated wheel
encoders, no slip --- so it is a floor the visual system must clear rather than a
replacement for it. It does, however, sharpen the thesis question: vision must
justify itself where dead reckoning fails, which means long runs, wheel slip,
relocalization after tracking loss, and moving obstacles.

\textbf{The detector's limit is detection, not semantics.} Over \num{9860}
frames the detector reaches precision 0.684 and moving-object recall 0.591, while
class assignment is 0.996 accurate once an object has been localized. Improving
the classifier would change nothing; recall is the binding constraint. The model
also produces bounding boxes rather than segmentation, and class membership is
not evidence of motion.

\textbf{The reconstructed map is geometrically sound.} Against the simulator's own
collision geometry at \SI{0.25}{\metre} tolerance, occupied-cell precision is
\SI{96.6}{\percent} with \SI{89.5}{\percent} recall over the region the robot
observed. The lower whole-building figure of \SI{67.2}{\percent} is route
coverage --- the run never entered the side aisles --- not mapping failure.

\textbf{Neither pipeline runs at input rate.} Frames arrive every
\SI{33}{\milli\second}; ORB-SLAM3 tracking takes \SI{63}{\milli\second} and its
local mapping \SI{190}{\milli\second}, and enabling the detector raises these to
\SI{113}{} and \SI{356}{\milli\second}. Input retention falls from
\SI{89}{\percent} to \SI{67}{\percent}.

\section{What is not established}

\textbf{The dynamic-scene results are withheld.} Those worlds are not yet a
physical simulation: props follow scripted kinematic paths, carrying collision
geometry but producing no dynamic response against the robot. Forty-eight runs
were logged there and all forty-eight fail the plausibility check, but those
numbers describe the simulation rather than the problem, so they are excluded
here and retained in the record as excluded evidence. The consequence is that the
condition this thesis is named for has \emph{no valid testbed} --- untested,
rather than tested and failed. Rebuilding those worlds is the first item of the
next period and the cheapest candidate explanation for the breakdown itself.

\textbf{The core hypothesis has no valid test behind it.} Two instrumentation
faults prevent it, and the environment question above compounds them. The VINS-Fusion filter never fired: all twelve runs requested
it and all twelve record zero matched detector frames and zero removed tracks,
even though the detector works offline. Separately, eleven of twelve ORB-SLAM3
baseline/filtered pairs received frame counts differing by \SIrange{10}{30}{\percent}
on the same recording. One admissible pair exists out of twenty-four. No
cross-treatment accuracy conclusion is drawn. Both are code-path faults and stand
regardless of scene, but the filtered runs were recorded in the dynamic worlds,
so the hypothesis now needs a working mask path \emph{and} a valid environment.

\section{How the work is validated}

Every comparison is tested for admissibility before its numbers are read: the
same recorded input to both treatments within \SI{5}{\percent}, the intended
treatment demonstrably active rather than merely requested, a common
ground-truth interval with the matched sample count reported, and physical
plausibility. Four claims were retracted or prevented this period by that
protocol --- a camera-rotation parsing bug found through map distortion while ATE
looked plausible throughout; a map wrongly judged unnavigable because its grid
was mis-seated; a drivability score discovered to have been measured against the
map's own construction run; and the two filter faults above. Each would otherwise
have entered the thesis as a result.

\section{Next period, in priority order}

\begin{enumerate}[leftmargin=2.2em]
  \item Repair the VINS mask delivery path. This alone unblocks the core
    hypothesis, and needs no additional data.
  \item Equalise ORB-SLAM3 input across paired runs so the filter comparison
    becomes admissible.
  \item Rebuild the dynamic worlds with physically simulated props, giving the
    titular condition a valid testbed.
  \item Add disturbance to the simulation --- wheel slip and encoder
    quantisation --- since both idealisations currently favour the
    dead-reckoning baseline.
  \item Record an independent route for map drivability, and extend the map
    evaluation to the large warehouse.
\end{enumerate}

\section{Guidance requested}

\begin{enumerate}[leftmargin=2.2em]
  \item \textbf{Scope.} Given that a wheel+IMU EKF currently beats both visual
    pipelines in static scenes, should the thesis argue for vision as the primary
    localization sensor, or reframe around the conditions under which vision adds
    value over dead reckoning?
  \item \textbf{Depth or breadth.} Rebuild the dynamic environment properly
    before measuring anything further in it, or extend to the real robot and the
    large warehouse with the static results as they stand?
  \item \textbf{Integration.} Relocalization, global optimization and
    multi-session map merging are listed as future work. Is that the right next
    contribution, or premature until dynamic-scene tracking is sound?
\end{enumerate}

\vspace{0.6em}
\noindent\rule{\textwidth}{0.4pt}\\[0.3em]
\footnotesize Full detail, including every retracted, excluded and failed
measurement, is in the accompanying progress report. Results flagged there in red
are measured but not yet trustworthy as evidence, and are excluded from the
claims above.

\end{document}
